% =====================================================================
%  Thème 4 — Produit scalaire et orthogonalité — DIFFICILE — Exercices
% =====================================================================
\documentclass[11pt,a4paper]{article}
\usepackage[utf8]{inputenc}
\usepackage[T1]{fontenc}
\usepackage{lmodern}
\usepackage[french]{babel}
\usepackage{amsmath,amssymb,mathtools}
\usepackage{xcolor}
\usepackage{tikz}
\usepackage{enumitem}
\usepackage{fancyhdr}
\usepackage[a4paper,margin=2.1cm,headheight=26pt,headsep=10pt]{geometry}
\usetikzlibrary{arrows.meta,calc}

\definecolor{primary}{RGB}{222,70,80}
\definecolor{primarydark}{RGB}{150,30,42}
\definecolor{vecu}{RGB}{20,20,20}
\definecolor{vecv}{RGB}{0,110,200}

\pagestyle{fancy}\fancyhf{}
\fancyhead[L]{\small\textcolor{primarydark}{\textbf{Fiche 2} — Calcul vectoriel}}
\fancyhead[R]{\small\textcolor{primarydark}{\textsc{Thème 4 — Difficile — Exercices}}}
\fancyfoot[C]{\small\thepage}
\renewcommand{\headrulewidth}{0.8pt}

\newcommand{\vc}[1]{\overrightarrow{#1}}
\providecommand{\norm}[1]{\left\lVert #1\right\rVert}
\newcommand{\exo}[1]{\medskip\noindent{\color{primary}\bfseries\large Exercice #1.}\par\nopagebreak\smallskip}

\begin{document}
\begin{center}
{\LARGE\bfseries\color{primarydark}Produit scalaire et orthogonalité — Niveau Difficile}\\[2pt]
{\color{primary}\rule{0.55\linewidth}{1.1pt}}
\end{center}
\vspace{1mm}
{\small\itshape Objectif : produit scalaire dans une figure, orthogonalité paramétrée, application au
travail d'une force, et synthèse « vrai / faux » de type concours.}
\vspace{2mm}

\exo{1} \textbf{Produit scalaire dans un carré (type concours).}
$PQRS$ est un carré de côté $2$, placé dans un repère orthonormé comme ci-dessous :
$P\,(0;2)$, $Q\,(2;2)$, $R\,(2;0)$, $S\,(0;0)$.

\begin{center}
\begin{tikzpicture}[scale=1.1]
  \coordinate (P) at (0,2); \coordinate (Q) at (2,2);
  \coordinate (R) at (2,0); \coordinate (S) at (0,0);
  \draw[black!70] (P)--(Q)--(R)--(S)--cycle;
  \draw[gray!55,dashed] (P)--(R); \draw[gray!55,dashed] (Q)--(S);
  \foreach \pt/\pos in {P/above left,Q/above right,R/below right,S/below left}{
    \fill (\pt) circle (1.6pt); \node[\pos,font=\small] at (\pt){$\pt$};}
\end{tikzpicture}
\end{center}

\begin{enumerate}[label=\textbf{\alph*)},leftmargin=2em,itemsep=3pt]
  \item Déterminer les composantes de $\vc{PQ}$ et de $\vc{RP}$.
  \item Calculer le produit scalaire $\vc{PQ}\cdot\vc{RP}$.
  \item Calculer $\vc{PR}\cdot\vc{QS}$ et en déduire une propriété des diagonales du carré.
\end{enumerate}

\exo{2} \textbf{Orthogonalité paramétrée.}
\begin{enumerate}[label=\textbf{\alph*)},leftmargin=2em,itemsep=3pt]
  \item Déterminer le réel $t$ pour lequel $\vec u\,(t+1;2)$ et $\vec v\,(2;t)$ sont orthogonaux.
  \item Déterminer toutes les valeurs de $t$ pour lesquelles $\vec a\,(t;3)$ et $\vec b\,(t;-12)$ sont
        orthogonaux.
\end{enumerate}

\exo{3} \textbf{Application — travail d'une force.}
Une force constante $\vec F\,(6;8)$ (en newtons) s'exerce sur un objet. Le travail de $\vec F$ lors
d'un déplacement $\vec d$ est $W=\vec F\cdot\vec d$ (en joules).
\begin{enumerate}[label=\textbf{\alph*)},leftmargin=2em,itemsep=3pt]
  \item Calculer la norme $\norm{\vec F}$ de la force.
  \item Calculer le travail pour le déplacement $\vec d_1\,(3;0)$.
  \item Calculer le travail pour le déplacement $\vec d_2\,(4;-3)$ et interpréter le résultat
        (relation entre $\vec d_2$ et $\vec F$).
\end{enumerate}

\exo{4} \textbf{Synthèse « vrai / faux ».}
Dans un repère orthonormé, on considère le carré $OHGF$ avec $O\,(0;0)$, $H\,(2;0)$, $G\,(2;2)$,
$F\,(0;2)$, et $E\,(0;1)$ le milieu de $[OF]$. Répondre par \textbf{Vrai} ou \textbf{Faux} en justifiant.
\begin{enumerate}[label=\textbf{\Alph*.},leftmargin=2.2em,itemsep=3pt]
  \item $\vc{OG}\cdot\vc{OH}=4$.
  \item Les vecteurs $\vc{OG}$ et $\vc{OH}$ sont orthogonaux.
  \item $\vc{OE}$ et $\vc{OH}$ sont orthogonaux.
  \item $\vc{OF}\cdot\bigl(3\,\vc{OF}\bigr)=12$.
\end{enumerate}

\end{document}
